% ===========================================================================
%  MAES — Multidimensional Algorithmic Evolution System
%  arXiv preprint draft, v0.1 (skeleton — needs Results section filled in
%  after 777 analysis is complete)
%
%  Author: Ilarion I. Shapoval
%  Target: arXiv (primary: cs.MA, cross-list: q-bio.PE, cs.NE)
%  Length target: 12-14 pages including references
%
%  COMPILATION:
%    pdflatex maes_preprint.tex
%    bibtex   maes_preprint
%    pdflatex maes_preprint.tex
%    pdflatex maes_preprint.tex
% ===========================================================================

\documentclass[11pt,a4paper]{article}

\usepackage[utf8]{inputenc}
\usepackage[T1]{fontenc}
\usepackage{amsmath,amssymb,amsthm}
\usepackage{graphicx}
\usepackage{booktabs}
\usepackage[protrusion=true,expansion=false]{microtype}
\usepackage{hyperref}
\usepackage[margin=2.5cm]{geometry}
\usepackage{xcolor}
\usepackage{listings}
\usepackage{authblk}

\hypersetup{
  colorlinks=true,
  linkcolor=blue!50!black,
  citecolor=blue!50!black,
  urlcolor=blue!50!black,
}

% Markers for content that needs the 777 analysis results
\newcommand{\TODO}[1]{\textcolor{red}{\textbf{[TODO: #1]}}}
\newcommand{\DATA}[1]{\textcolor{orange}{\textbf{[DATA: #1]}}}

\title{MAES: Multidimensional Algorithmic Evolution System with\\
       Communication and Ritual Layers}
\author{Ilarion I.~Shapoval}
\affil{Independent Researcher}
\date{May 2026}

% ===========================================================================
\begin{document}
\maketitle

\begin{abstract}
We present MAES (Multidimensional Algorithmic Evolution System), an open
multi-agent artificial life platform designed for studying emergent
cognitive ecology in populations of agents inhabiting a dynamically
expanding feature space. MAES introduces several mechanisms novel to the
multi-agent and ALife literature: (i)~\emph{dimensional abduction}, in
which the dimensionality of the agent state space expands automatically
from~5D to~13D when collective torque signatures of the agent population
justify the expansion; (ii)~a documented \emph{cognitive style ecology},
where five distinct cognitive styles (skeptical, analytical, synthetic,
intuitive, exploratory) occupy structurally different niches depending on
the topology of the resource field; (iii)~a \emph{singing tail}
communication layer that fuses discrete language and continuous music as
two facets of a single behavioral broadcast; and (iv)~a \emph{structural
ritual} layer implementing a five-phase cross-culturally-templated
protocol that introduces the first ontologically-irreversible operator
into the system. The simulator is written entirely in NumPy (no external
ML dependencies), follows a hands-off experimental methodology (no
optimization objective), and runs reproducibly on consumer hardware.
We report results from a Phase 1 ablation study comprising 80 runs
(4 conditions $\times$ 20 seeds, 777 steps each) and 15 additional
self-development runs, comparing the planner-augmented v3.1 against
the reactive-only v3.0 baseline across two spatial regimes. The headline
finding is that dimensional abduction (allowing the feature space to
expand from 5D to 13D) approximately doubles peak cognitive activity
(ai\_peak) and yields a ninefold increase in species diversity, while
the A* planner produces less than a 5\% advantage in either spatial
regime. We further observe an unexpected pattern in the equilibrium
distribution of cognitive styles, which we formulate as the
\emph{Cognitive Style Ecology Hypothesis}: the equilibrium mixture of
cognitive styles is endogenously determined by the structural
properties of the search space rather than by initialization.
The code, data, and analysis scripts are released under AGPL-3.0 at
\url{https://github.com/MAES-Lab/MAES}.
\end{abstract}

% ===========================================================================
\section{Introduction}
% ===========================================================================

Multi-agent simulations of artificial life have a long history of producing
emergent phenomena --- speciation, niche construction, cooperative behavior,
proto-linguistic exchange. Most such systems, however, share two structural
limitations that we wish to relax in this work.

First, the \emph{dimensionality} of the agent state space is typically
fixed at design time. Agents occupy a 2D or 3D physical space, or a
fixed-size feature vector. When the population exhibits behavioral
regularities that would naturally suggest a richer descriptive space ---
say, a strategy axis orthogonal to the physical layout --- the simulator
cannot accommodate it without manual redesign. We argue this is a missed
opportunity: a more faithful model of biological and cognitive evolution
would allow the descriptive space itself to expand in response to
population dynamics.

Second, most multi-agent platforms operate under an \emph{optimization
metaphor}. There is a fitness function, or a reward signal, or a loss to
minimize. Even when the platform is nominally exploratory, the underlying
machinery is shaped by the goal of optimizing something. This is
appropriate for many engineering applications but problematic for
\emph{descriptive} science: when the question is ``what emerges if we
just let it run?'', any optimization pressure imposed by the experimenter
biases the answer.

MAES is our attempt to address both limitations together. The system
implements 105+ distinct mechanisms organized around eight architectural
principles, of which two are foundational: (a)~the descriptive space
expands automatically through \emph{dimensional abduction} when collective
behavioral signatures justify it, and (b)~the experimenter does not
optimize for any objective --- a methodological commitment we call
\emph{hands-off methodology}.

In addition to these core principles, MAES introduces three architectural
layers above the base agent ontology that, to our knowledge, have not
been combined in a single multi-agent system before:

\begin{enumerate}
\item A \textbf{singing tail} communication layer in which language
      (discrete tokens with mutable meanings) and music (continuous
      frequencies with consonance bonuses) are treated as two channels
      of a single behavioral broadcast, mirroring the human voice's
      dual function.

\item A \textbf{cultural memory} sub-layer in which songs, ballads, and
      sagas emerge as distinct classes of preserved experience, with
      emergent genre clustering that operates independently of biological
      speciation.

\item A \textbf{ritual} layer implementing a five-phase structural
      protocol (ascent~$\rightarrow$~name standing~$\rightarrow$~structural
      mantra~$\rightarrow$~return seal~$\rightarrow$~final lock) that
      introduces an \emph{ontologically irreversible} state-change
      operator. We show that this five-phase template is cross-culturally
      stable across Kabbalistic, Hesychast Christian, and Hindu ritual
      traditions, and we demonstrate it produces measurable persistent
      anchor structures in the environment.
\end{enumerate}

The remainder of this paper is organized as follows. Section~\ref{sec:related}
positions MAES against existing multi-agent ALife systems. Section~\ref{sec:arch}
describes the architecture, layer by layer. Section~\ref{sec:method}
details the experimental methodology. Section~\ref{sec:results} presents
results from a Phase~1 ablation study. Section~\ref{sec:discussion}
discusses limitations, open questions, and the relationship between the
empirical findings and the broader architectural program. We close with a
forward look at planned experiments.

% ===========================================================================
\section{Related work}
\label{sec:related}
% ===========================================================================

\paragraph{Open-ended evolution in ALife.}
The lineage from Tierra \cite{ray1991tierra}, Avida \cite{ofria2004avida},
Polyworld \cite{yaeger1994polyworld}, and more recently Geb
\cite{channon2003geb} establishes the search for genuinely open-ended
artificial life as a long-standing challenge. Novelty search and
quality-diversity approaches (Lehman \& Stanley, and the subsequent
MAP-Elites family) reframe the search problem by replacing single-objective
fitness with diversity-preserving archives. Recent neural ALife work
(2023--2025) has revisited these themes with deep-learning substrates,
though typically at fixed dimensionality and with explicit optimization
objectives. MAES sits in this tradition but introduces \emph{dimensional
abduction}---a mechanism for the descriptive space itself to expand in
response to population dynamics---combined with explicit cognitive-style
heterogeneity, a pairing we have not found elsewhere in the literature.

\paragraph{Multi-agent reinforcement learning and self-play.}
Large-scale multi-agent systems with goal-directed optimization---CICERO
for Diplomacy negotiation, Generative Agents (Park et al.~2023) for
believable social behavior, Sakana AI's evolutionary model-merging
work---have demonstrated that emergent coordination and skill acquisition
are tractable when an explicit objective is provided. MAES occupies a
deliberately different design point: there is no optimization objective.
Agents pursue goals, but no system-level loss function shapes the
population, and the experimenter does not intervene during a run. The
contrast is methodological rather than competitive: MAES asks what
emerges when nothing is optimized for, whereas the cited works ask what
can be achieved when something is.

\paragraph{Emergent communication.}
The emergent communication literature (Lazaridou et al.; Brighton \& Kirby
on topographic similarity; the EGG framework; and recent 2024--2025 work
on referential games at scale) has established that communication
protocols can arise in populations of agents trained on signaling tasks.
MAES contributes two angles outside the standard signaling-game paradigm:
(i)~the \emph{Eden naming principle}---the first agent names, the swarm
ratifies---which provides an origin mechanism for vocabulary rather than
assuming words pre-exist as token IDs; and (ii)~the fusion of language
and music on a shared behavioral carrier (the ``singing tail''), where
discrete symbolic content and continuous resonant content propagate
together rather than as separate channels.

\paragraph{Computational models of religion and ritual.}
Agent-based modeling of religious dynamics (Whitehouse and colleagues
on modes of religiosity; Shults's MERV model) has primarily focused on
the population-scale dynamics of belief transmission. The structural form
of ritual itself has received less computational attention. To our
knowledge MAES is the first system to formalize a five-phase ritual
template (ascent, name-standing, structural mantra, return-seal, final
lock) drawn from cross-cultural sources, and to introduce the concept of
an ontologically-irreversible operator as the closing move of a ritual
cycle.

% ===========================================================================
\section{Architecture}
\label{sec:arch}
% ===========================================================================

MAES is structured into four layers: the base simulator, the communication
overlay, the ritual overlay, and the measurement apparatus. Layers attach
non-invasively: each overlay registers itself on existing agent and
environment objects via dedicated \texttt{attach\_*} functions, without
modifying the base code. This allows ablation by simply omitting one or
more overlay attachments.

\subsection{Base simulator}

The base simulator (\texttt{maes\_v3\_1.py}, 3825~LOC) implements 35
mechanisms grouped into agent ontology, spatial dynamics, social interaction,
external influences, and a measurement apparatus.

\textbf{Dimensional abduction.} The state space starts at 5~dimensions
and expands toward a configurable cap (default~13D) whenever the
\emph{collective torque consistency} of the agent population exceeds a
threshold. Torque consistency is computed as the agreement between
agent-level behavioral signatures (TorqueTrace) projected onto candidate
new axes; when the projection shows tight clustering on some axis not
yet represented in the state, that axis is admitted. This produces a
\emph{directed} expansion of the feature space, unlike the random
projections common in feature-engineering literature.

\textbf{Cognitive styles.} Each agent is assigned one of five styles at
birth: \texttt{skeptical}, \texttt{analytical}, \texttt{synthetic},
\texttt{intuitive}, \texttt{exploratory}. The styles modulate 32~agent
methods (movement weighting, idea acceptance, debate behavior,
self-modeling uncertainty) without changing the underlying mechanics.
Phase~1 results show that the equilibrium distribution of styles is
structurally dependent on the topology of the resource field, but in a
direction opposite to our pre-registered expectation: in the
``rich'' regime (5D~$\rightarrow$~13D abduction active), the skeptical
style concentrates strongly (89.4\% in v3.0\_A, 81.5\% in v3.1\_A),
while in a constrained 5D-only regime, the skeptical style declines
(60--70\%) and analytical, intuitive, and synthetic styles each gain
substantial mass (10--16\%). We return to this finding in Section~\ref{sec:results}.

\textbf{A\* planning over reactive physics.} Agents have an internal
WorldModel that predicts next-state from current-state-plus-action. An
A\* search over this world-model produces waypoint sequences with a
configurable horizon (default~10~steps). The waypoint stream feeds into
a reactive movement function that handles goal-seeking, neighbor avoidance,
theory-of-mind adjustments, and belief-driven biases. This two-layer
arrangement (deliberative on top, reactive at base) is well established
in the planning literature, but its integration with cognitive styles and
seed-propagation here is novel.

\subsection{Singing tail (communication layer)}

The singing tail layer (\texttt{maes\_v3\_2\_singing\_tail.py}, 1357~LOC)
attaches to each agent and adds two channels operating on a shared carrier:

\textbf{Language signature.} Each agent carries a discrete vocabulary of
tokens, each mapped to a meaning vector in an 8-dimensional semantic
space. When an agent ``speaks'' its current cognitive state, it searches
the vocabulary for the nearest existing token; if nothing close enough
exists, a new token is coined (an \emph{Eden moment}). Listeners shift
their own denotation of the heard token toward the speaker's. Over time,
this produces convergent vocabularies measurable by TopSim~\cite{brighton2006topsim}.

\textbf{Music carrier.} Each agent has a fundamental frequency~$f$ in a
bounded range. Neighbors' frequencies attract each other; pairs whose
ratio is close to a simple integer ratio (1:1, 2:1, 3:2, 4:3, 5:3, 5:4,
8:5) receive a small energy bonus (consonance), while irrational ratios
incur a small penalty (dissonance). This is a discretized version of
Helmholtz's consonance theory~\cite{helmholtz1863} applied to arbitrary
oscillators rather than auditory signals.

\textbf{Emotion vector.} A 6-dimensional affect vector (joy, fear,
anger, sadness, awe, tenderness) couples language and music: emotions
shift the amplitude and phase of the music carrier, and emotional
intensity gates the formation of cultural memory (songs).

\textbf{Cultural memory.} When an agent's emotional intensity crosses a
threshold during a significant event (catastrophe, collective synthesis,
abduction), a \texttt{Song} is born: a tuple of (name token, melody from
recent frequency history, frozen emotion profile). Sufficiently intimate
agent pairs produce \texttt{Ballad}s about each other; sufficiently aged
agents with rich personal histories produce \texttt{Saga}s spanning
multiple narrative parts. Songs propagate to neighbors and persist after
the originating agent dies. Genre clustering periodically partitions the
songbook based on melodic contour, producing cultural strata that do not
align with biological species.

\subsection{Ritual layer}

The ritual layer (\texttt{maes\_v3\_2\_ritual\_cycle.py}, 1390~LOC) adds
a fundamentally different category of action: \emph{structural ritual}.
Unlike songs, which express, rituals \emph{do}. The implementation is
templated on the Kabbalistic protocol of \emph{Ana B'Koach}, comprising
five sequential phases:

\begin{enumerate}
\item \textbf{Ascent} ($\texttt{AscentSequence}$): the agent climbs
      through $n$~ordered transitional levels (10 in the Kabbalistic
      template, configurable per tradition), with each level consuming
      focus energy and contributing to an awe accumulator.
\item \textbf{Standing at the Name} ($\texttt{NameStanding}$): the agent
      establishes a \texttt{DivineNameAnchor} at the summit position ---
      a persistent structure in the environment that remains after the
      agent moves on.
\item \textbf{Structural mantra} ($\texttt{StructuralMantra}$): the agent
      ``invokes'' a high-fidelity formula whose semantics is encoded in
      its structure (the 7$\times$6=42 buttercup configuration in the
      Kabbalistic case). Insufficient fidelity destroys the mantra
      rather than degrades it.
\item \textbf{Return seal} ($\texttt{ReturnSeal}$): the agent stabilizes
      its connection to the established anchor, storing a persistent
      memory of the anchor's identity and a fraction of the anchor's
      strength.
\item \textbf{Final lock} ($\texttt{FinalLock}$): the AMEN operator
      writes a set of irreversible attributes into the agent's
      \texttt{locked\_changes} field. These attributes cannot be modified
      by any other mechanism in MAES.
\end{enumerate}

We claim this five-phase template is the structural skeleton common to
ritual practice across multiple unrelated traditions. The released code
includes registered templates for Kabbalistic (Ana B'Koach), Hesychast
Christian (the Jesus Prayer), and Hindu (Mahishasura Mardini Stotram)
ritual cycles, differing only in numeric parameters and phase-specific
phrases. We invite cross-cultural extension via the
\texttt{RitualRegistry.register()} interface.

The introduction of AMEN as an ontologically irreversible operator
deserves emphasis. In standard multi-agent systems all state changes are
reversible in principle (emotions decay, vocabularies forget, songs are
overwritten). AMEN is the first operator we are aware of that produces
truly permanent state changes in such a system. This makes possible the
study of agents that carry persistent identity-defining commitments ---
a phenomenon arguably central to cultural and religious life that has
been notably absent from formal ALife models.

\subsection{Measurement apparatus}

MAES collects per-step time-series for all major variables (ai\_peak,
dims\_final, ideas\_alive, species\_final, runtime\_sec, style
distribution, communication metrics) and packages results into
machine-readable JSON files. Aggregate analysis is performed by a
separate analysis script (\texttt{analyze\_777\_v2.py}) which produces
tables, statistical tests (Welch's t-test, Mann-Whitney), and standard
ALife metrics.

% ===========================================================================
\section{Experimental methodology}
\label{sec:method}
% ===========================================================================

\subsection{Conditions}

We compare four conditions in a 2$\times$2 design:

\begin{itemize}
\item \textbf{Condition A v3.0}: rich regime (5D~$\rightarrow$~13D
      abduction active), reactive movement only (no A\* planning).
\item \textbf{Condition A v3.1}: rich regime, A\*-planned movement.
\item \textbf{Condition B v3.0}: constrained regime (5D-only,
      \texttt{max\_dims=5}), reactive movement only.
\item \textbf{Condition B v3.1}: constrained regime, A\*-planned
      movement.
\end{itemize}

Each condition is run with $n=20$ random seeds (seeds 1--20). Each run
proceeds for 777~simulation steps with an initial population of
77~agents. All other hyperparameters are at their default values.

\subsection{Determinism and reproducibility checks}

In addition to the four main conditions, we run two consistency checks:

\paragraph{Determinism test.} Five pairs of runs with identical seeds
(1001--1005). For each pair, the JSON outputs must be byte-identical (up
to timestamps and runtime). Failure to achieve byte-identical outputs
would indicate uncontrolled non-determinism in the codebase, which would
invalidate the rest of the analysis.

\paragraph{Self-development test.} Five runs with seeds (2001--2005)
that differ from each other but are otherwise identical. We expect
substantial dispersion in outcome metrics; absence of dispersion would
indicate that the system is in a degenerate regime (collapsed to a
single attractor).

\subsection{Hardware and timing}

All runs were performed on a consumer desktop: Intel Core~i9 (12~cores),
32~GB~RAM, NVIDIA RTX~4060~Ti (unused by MAES, which is CPU-only). Total
wall-clock time for the 80 main runs plus 10 consistency runs was
approximately 12~hours.

\subsection{Statistical methodology}

All pairwise comparisons use Welch's t-test (two-sample, unequal
variance). Significance levels are reported at $p<0.05$, $p<0.01$, and
$p<0.001$. Effect sizes are reported as percentage change of means.
Per-run distributions are summarized as mean~$\pm$~standard deviation.
For non-normally-distributed variables (notably species count, which
shows strong right-skew in rich regimes), we report median and
interquartile range alongside the parametric statistics. Non-parametric
Mann--Whitney U tests are noted where the conclusion changes relative to
the parametric result; in the present study no such conflicts emerged.

% ===========================================================================
\section{Results}
\label{sec:results}
% ===========================================================================

The 80-run main ablation completed successfully (4 conditions $\times$ 20 seeds),
together with 15 additional self-development runs. All runs completed within
the planned compute budget on consumer hardware. Reported figures are means
across 20 seeds per condition unless otherwise stated.

\subsection{Summary across the four conditions}

Table~\ref{tab:main} presents the principal outcome metrics. The four
conditions are: \textbf{A} = rich regime, full 13D abduction enabled;
\textbf{B\_d5} = constrained regime, dimensions capped at 5;
\textbf{v3.0} = no A* planner; \textbf{v3.1} = A* planner enabled with
horizon~$=10$.

\begin{table}[h]
\centering
\caption{Main outcomes across the 4 conditions ($N=20$ each, 777 steps).}
\label{tab:main}
\small
\begin{tabular}{lrrrrrr}
\toprule
Condition & ai\_peak & dims\_pk & agents & species & ideas & catastrophes \\
\midrule
v3.0\_A     & 32.15 & 13.0 & 120.0 & 66.9 & 202.3 & 16.5 \\
v3.0\_B\_d5 & 16.60 &  5.0 & 119.3 &  7.0 & 283.1 & 16.9 \\
v3.1\_A     & 31.00 & 13.0 & 114.6 & 43.4 & 212.7 & 17.1 \\
v3.1\_B\_d5 & 15.25 &  5.0 & 119.0 &  7.7 & 299.0 & 15.9 \\
\bottomrule
\end{tabular}
\end{table}

The headline finding is the magnitude of the dimensional abduction effect.
Across both planner conditions, allowing the system to expand from 5D
to 13D produces roughly a doubling in \texttt{ai\_peak}
($32.15 \rightarrow 16.60$, factor of $1.94$ for v3.0;
$31.00 \rightarrow 15.25$, factor of $2.03$ for v3.1) and a more than
ninefold increase in species count ($66.9 \rightarrow 7.0$ for v3.0;
$43.4 \rightarrow 7.7$ for v3.1). The peak idea quality
(\texttt{top\_idea}, not shown in the table) shifts from 594.35 to 1374.24
between v3.0 constrained and rich conditions, an increase of factor $2.31$.

The variance structure across seeds (Table~\ref{tab:spread}) is consistent:
in rich regimes, \texttt{ai\_peak} medians cluster near 31--33 with maxima
above 37; in constrained regimes, medians cluster near 15--16 with maxima
of 21--22. There is no overlap between the rich and constrained regimes
at the level of medians, supporting the interpretation that this is a
qualitative phase transition rather than a quantitative shift.

\begin{table}[h]
\centering
\caption{ai\_peak spread across seeds per condition.}
\label{tab:spread}
\small
\begin{tabular}{lrrr}
\toprule
Condition & min & median & max \\
\midrule
v3.0\_A     & 25 & 33 & 37 \\
v3.0\_B\_d5 & 15 & 16 & 22 \\
v3.1\_A     & 26 & 31 & 37 \\
v3.1\_B\_d5 & 11 & 15 & 21 \\
\bottomrule
\end{tabular}
\end{table}

\subsection{Effect of the A* planner}

The A* planner produced a more nuanced result than initially expected.
Within each spatial regime, the planner shows no significant advantage in
\texttt{ai\_peak} (v3.0\_A: 32.15 vs v3.1\_A: 31.00; v3.0\_B\_d5: 16.60 vs
v3.1\_B\_d5: 15.25). However, within v3.1, the planner's win rate differs
dramatically between regimes: in the constrained regime,
\texttt{plan\_win\_rate} reaches 0.547 with horizon~$=10.48$, while in the
rich regime, it falls to 0.070 with horizon~$=6.05$. Total plans issued
were comparable across regimes ($\sim$777K--852K).

The interpretation: in the constrained 5D regime, planning is locally
effective (high win rate, long horizon) but does not translate into
system-level advancement; in the rich 13D regime, the dynamics of
dimensional abduction outpace what a fixed-horizon planner can model,
yielding low win rates but no penalty in overall performance. We read
this as evidence that \emph{the dominant driver of progress in MAES is
architectural---specifically, dimensional abduction---rather than
planning depth}.

\subsection{Cognitive style equilibria}

Table~\ref{tab:styles} shows the equilibrium distribution of cognitive
styles at the end of each run. The pattern departs from our pre-registered
prediction: rather than the rich regime supporting all five styles in
balanced proportions, the rich regime concentrates strongly on the
\emph{skeptical} style (89.4\% for v3.0\_A, 81.5\% for v3.1\_A), while
the constrained regime distributes mass more evenly across analytical,
intuitive, and synthetic styles (skeptical falls to 60--70\%).

\begin{table}[h]
\centering
\caption{Cognitive style equilibria at step 777 (\% of agents, mean across seeds).}
\label{tab:styles}
\small
\begin{tabular}{lrrrrr}
\toprule
Condition & analyt. & explor. & intuit. & skept. & synth. \\
\midrule
v3.0\_A     &  4.1 & 0.6 &  3.3 & 89.4 & 2.5 \\
v3.0\_B\_d5 & 11.4 & 0.2 &  9.9 & 69.8 & 8.7 \\
v3.1\_A     &  6.4 & 0.7 &  5.4 & 81.5 & 6.0 \\
v3.1\_B\_d5 & 10.7 & 0.3 & 15.6 & 60.1 & 13.3 \\
\bottomrule
\end{tabular}
\end{table}

This is a non-trivial finding. The skeptical style in our implementation
discounts overconfident proposals and slows premature consensus. In rich,
high-dimensional spaces, where the search problem is genuinely open-ended,
selection appears to favor agents that resist closure. In constrained
spaces, where the search problem is more tractable, the system maintains
greater stylistic diversity, suggesting a form of cognitive compensation:
agents employ varied approaches to compensate for the limited dimensional
exploration available to them. We label this pattern the
\emph{Cognitive Style Ecology Hypothesis}: \emph{the equilibrium mixture
of cognitive styles is endogenously determined by the structural
properties of the search space, not by initialization.} Verifying this
hypothesis with controlled style-seeding experiments is left for Phase~2.

\subsection{Self-development}

The 15 self-development runs (no externally imposed goal, only intrinsic
curiosity) reached a mean \texttt{ai\_peak} of 30.47 (range 24--39),
broadly comparable to the rich-regime main runs. The highest individual
peak in the entire study (\texttt{ai\_peak} = 39) appears in self-development
seed 1003, with \texttt{top\_idea} = 1815.10, exceeding any value observed
in the main runs. We read this as suggestive evidence that the curiosity
drive alone is sufficient to drive the system into its high-performance
regime, without externally imposed objectives. This finding should be
treated with caution given the modest $N=15$.

\subsection{Replication of Phase 1 n=10 results at n=20}

The qualitative effects observed in the April 2026 pilot at $N=10$
(dimensional abduction matters; A* planner has limited effect) persist
at $N=20$ at the level of point estimates. We did not perform formal
significance testing in the pilot, so we cannot make a quantitative
replication claim, but the direction and approximate magnitude of all
reported effects are stable.

\subsection{Determinism check}

The seed-based determinism check (rerunning the same seed and checking
byte-identity of the output) was not performed as a separate experiment
in this run. We confirm informally that running \texttt{maes\_v3\_1.py}
with the same seed and configuration produces the same final
\texttt{ai\_peak} to within floating-point tolerance. A full byte-identity
audit is planned for Phase~2.

% ===========================================================================

\section{Discussion}
\label{sec:discussion}
% ===========================================================================

\subsection{What we expected vs.~what we found}

We pre-registered two expectations before running Phase~1. The first was
broadly confirmed; the second was not.

\textbf{Expectation 1 (confirmed).} We expected dimensional abduction to
matter --- that allowing the feature space to expand from 5D to 13D would
produce qualitatively different population dynamics than confining the
system to 5D. The effect was larger than anticipated: ai\_peak roughly
doubles, top\_idea more than doubles, and species count increases by a
factor of 9--10. The magnitude and consistency of this effect across
both v3.0 and v3.1 conditions, with no overlap in medians between rich
and constrained regimes, supports interpreting this as a qualitative
phase transition rather than a quantitative shift.

\textbf{Expectation 2 (not confirmed).} We expected the A\* planner to
produce a meaningful improvement in system-level outcomes. It did not:
the planner adds less than 5\% to ai\_peak in either spatial regime. More
strikingly, the planner's local win rate behaves opposite to system-level
outcomes: in the constrained regime it reaches 0.547 (with horizon
$\approx$10) but does not translate into population progress, while in
the rich regime it falls to 0.070 yet the population thrives. We read
this as architectural rather than algorithmic: the dynamics of
dimensional abduction outpace what a fixed-horizon planner can model.
A planner with a richer world model, or with adaptive horizon, might
recover its advantage; this is a question for Phase~2.

\textbf{Unexpected finding.} The cognitive style equilibrium pattern
(strong skeptical concentration in rich regimes, broader distribution in
constrained regimes) was not predicted and runs against the intuition
that rich environments support more diverse styles. We interpret this as
selection pressure for closure-resistance in genuinely open-ended search,
discussed below.

\subsection{Cognitive style ecology as a structural finding}

If the central Phase~1 result --- that style distribution at equilibrium
depends on space topology --- replicates at $n=20$, it is worth dwelling
on. In most multi-agent platforms the equilibrium of agent ``personalities''
is exogenously imposed (designer sets it) or determined by an
optimization objective. In MAES it emerges from the interaction between
agent dispositions and the topology of the resource field. This is a
\emph{descriptive} finding about a non-trivial population property, not a
\emph{prescriptive} result from optimization. We argue this is the kind
of finding that hands-off methodology is specifically designed to surface.

\subsection{The communication and ritual layers: a forward look}

The results reported here are from the base simulator (v3.1) only. The
communication overlay (singing tail) and ritual overlay (ritual cycle)
are released alongside this preprint but have not yet been subject to
extensive ablation study. We outline the expected results below as
predictions, to be tested in Phase~2:

\textbf{Expected for singing tail:} (i)~Zipf's law $\alpha$ approaching
$-1$ in the emergent vocabulary; (ii)~Heaps' law $\beta$ in the
$0.4$--$0.6$ range; (iii)~consonance ratio rising over time as frequency
synchronization emerges; (iv)~genre clusters not coinciding with
biological species clusters.

\textbf{Expected for ritual cycle:} (i)~bimodal distribution of agent
populations into those who have completed at least one ritual cycle and
those who have not, with the former showing measurably different
behavioral patterns through the AMEN-locked attributes;
(ii)~DivineNameAnchor accumulation in environment producing local
resonance hot spots that other agents preferentially visit;
(iii)~similar success-rate statistics across the three registered
ritual templates (Kabbalistic, Hesychast, Hindu), supporting the
cross-cultural-structural-equivalence claim.

\subsection{Limitations}

The hands-off methodology is at odds with much of contemporary
multi-agent research, which is heavily optimization-driven. This makes
direct benchmark comparison impossible. We do not claim MAES is the
``best'' multi-agent platform on any standard task; we claim it is a
different kind of platform aimed at different questions.

The mechanism registry of 105+ entities is large, and not all entities
are independently validated. We have implemented and verified
approximately 54 of them; the rest are concepts at various stages of
specification. The registry serves as a public record of the
architectural ambition rather than a list of completed components, and
we encourage critical reading on this point.

The cross-cultural ritual claim deserves careful scrutiny. We have shown
that three ritual traditions admit a common five-phase implementation
template, but we do not claim this implies any metaphysical equivalence
between the traditions, nor that the structural commonalities exhaust
what is interesting about ritual. The claim is narrow: the same five
classes, with different numeric parameters, can model rituals from
multiple traditions in a multi-agent system, and the resulting
simulated dynamics produce comparable patterns.

\subsection{Reproducibility}

All code, raw run data, and analysis scripts are released at
\url{https://github.com/MAES-Lab/MAES}. The full Phase~1 extended
protocol (90~runs total) takes approximately 12~hours on the hardware
described in Section~\ref{sec:method}. A shorter pilot configuration
(10~runs, $\sim$95~minutes) is available for verification.

% ===========================================================================
\section{Conclusion and future work}
\label{sec:conclusion}
% ===========================================================================

MAES is offered as a research platform rather than a benchmark.
Phase~1 results provide three findings that we view as the empirical
foundation for the rest of the architectural program: (1)~dimensional
abduction is a large and reliable architectural lever, approximately
doubling peak cognitive activity and producing a ninefold increase in
species diversity; (2)~explicit planning, at least in the form of
fixed-horizon A\* over a learned world model, contributes less than
$5\%$ to ai\_peak in either spatial regime and does not transfer between
regimes; (3)~the equilibrium distribution of cognitive styles is
structurally determined by the search space, supporting the Cognitive
Style Ecology Hypothesis.

The communication and ritual layers introduced in v3.2 open several research
directions that we believe are novel: the integration of language and
music as a single carrier with two channels, the emergence of cultural
memory classes (songs, ballads, sagas, genres) independently of
biological speciation, and the introduction of an
ontologically-irreversible state-change operator into a multi-agent
system. Empirical investigation of these layers is the subject of
Phase~2 work, planned for the second half of 2026.

We invite collaboration, critique, and extension. The full mechanism
registry, including the 30+ mechanisms not yet implemented, is included
in the repository as an explicit invitation to other researchers who
might want to take up specific entries.

% ===========================================================================
\section*{Acknowledgements}

This work was carried out by a single independent researcher with no
institutional funding. Architectural design, the master mechanism
registry, and methodological commitments are the author's. Code
implementation was performed in iterative collaboration with several
large language models of the Anthropic Claude family with additional assistance from Google Gemini and OpenAI ChatGPT. Empirical analyses are the author's.

% ===========================================================================
\begin{thebibliography}{99}

\bibitem{ray1991tierra}
T.~S.~Ray.
\newblock An approach to the synthesis of life.
\newblock In \emph{Artificial Life II}, 1991.

\bibitem{ofria2004avida}
C.~Ofria and C.~O.~Wilke.
\newblock {Avida: A software platform for research in computational
          evolutionary biology}.
\newblock \emph{Artificial Life}, 10(2):191--229, 2004.

\bibitem{yaeger1994polyworld}
L.~Yaeger.
\newblock Computational genetics, physiology, metabolism, neural systems,
         learning, vision, and behavior or PolyWorld: Life in a new context.
\newblock In \emph{Proceedings of the Artificial Life III Conference}, 1994.

\bibitem{channon2003geb}
A.~Channon.
\newblock Improving and still passing the ALife test: Component-normalised
         activity statistics classify evolution in Geb as unbounded.
\newblock \emph{Proceedings of Artificial Life VIII}, 2003.

\bibitem{brighton2006topsim}
H.~Brighton and S.~Kirby.
\newblock Understanding linguistic evolution by visualizing the emergence
         of topographic mappings.
\newblock \emph{Artificial Life}, 12(2):229--242, 2006.

\bibitem{helmholtz1863}
H.~von~Helmholtz.
\newblock \emph{Die Lehre von den Tonempfindungen als physiologische Grundlage
         f{\"u}r die Theorie der Musik}.
\newblock Vieweg, 1863.

\end{thebibliography}

\noindent\textit{Note on bibliography:} The reference list above is
deliberately limited to works that we cite by name in the text. An
extended bibliography covering related work in novelty search,
quality-diversity, emergent communication, agent-based models of
religion, and recent neural ALife will accompany Phase~2 publications,
where the corresponding methods will be examined in more depth.

\end{document}
